\documentclass[10pt,a4paper]{amsart}
\usepackage[margin=19mm]{geometry}
\usepackage{amsmath,amssymb,mathtools}
\usepackage[scr=rsfs,cal=euler]{mathalfa}
\usepackage{xfrac}
\usepackage[unicode,hidelinks,bookmarks=false]{hyperref}
\newcommand{\ie}{i.e.\ }
\newcommand{\cf}{cf.\ }
\newcommand{\resp}{respectively\ }
\newcommand{\Subsection}{\subsection}
\providecommand{\nobreakdash}{\nobreak\hbox{-}}
\setlength{\emergencystretch}{2em}
\multlinegap0pt
\usepackage{bm}
\usepackage{bbm}
\usepackage{dsfont}
\usepackage{xspace}
\usepackage{cancel}
\usepackage{mathtools}
\usepackage{amsthm}
\makeatletter
\@ifundefined{definition}{\newtheorem{definition}{Definition}}{}
\@ifundefined{Lemma}{\newtheorem{Lemma}[definition]{Lemma}}{}
\makeatother
\newcommand{\C}{\mathbb{C}} % komplexe Zahlen
\newcommand{\E}{\mathbb{E}} % EW
\newcommand{\R}{\mathbb{R}} % reelle Zahlen
\title[MA(1) processes with uniform innovations conditioned to stay]{MA(1) processes with uniform innovations conditioned to stay positive in the non-expanding regime}
\author{Frank Aurzada, Virginia Worf}
\date{Source: arXiv:2606.01832v1 (CC BY 4.0)}
\begin{document}
\maketitle

\noindent\emph{Source and licence.} MA(1) processes with uniform innovations conditioned to stay positive in the non-expanding regime. Version arXiv:2606.01832v1 (CC BY 4.0), distributed under the Creative Commons Attribution 4.0 International licence, as recorded in the arXiv licence metadata for this exact version and shown on the abstract page https://arxiv.org/abs/2606.01832v1. Full rights notice: licenses/arxiv-2606.01832-CC-BY-4.0.txt.\par\smallskip

\noindent\textit{Editorial scope.} Counted unit: the Lemma whose source label is \texttt{asymptotic\_gelb\_eigenfunction} in the article (source0.tex lines 989--1003, source label \texttt{asymptotic\_gelb\_eigenfunction}), delivered together with the article's own complete original proof of it (source0.tex lines 1004--1139, one \texttt{proof} environment of 136 lines). No mathematics is written, repaired or re-derived: statement and proof are the article's own text line for line. Required context retained uncounted: asymptotic\_powerseries (lines 535--544); root\_blue (lines 575--577); asymptotic\_blau\_erste (lines 616--622); generatingfunction\_gelbe\_region (lines 888--898); root\_yellow\_2 (lines 939--941); lemma\_root\_yellow (lines 945--947). Every definition, display and label of those lines is retained unchanged. Omitted from this package and not redistributed: 1--534, 545--574, 578--615, 623--887, 899--938, 942--944, 948--988, 1140--1940. Nothing inside a retained excerpt is omitted or altered: every retained body is delivered exactly as the article's lines are written, so no omission receipt of any kind accompanies this package. Citations: the retained excerpts carry no citation command at all, so no bibliography entry of the article is transcribed and no reference is redistributed. Reference adaptation: none was needed and none was made; every reference inside the delivered unit resolves inside it, so no unresolved reference or citation is produced. Printed slips kept literal: no printed word, symbol, hypothesis or number of the counted statement or of its proof was altered. Layout and class: the article's own class is \texttt{article} with packages including babel, todonotes, epsfig, dsfont, comment, subfig, amsmath, amssymb, amsthm, thmtools, tikz, xparse, ifthen, textcomp; the wrapper is the lane's pinned \texttt{amsart} editorial wrapper at 10pt on A4, which restates the theorem-like environments the retained text needs and adds the article's own macro definitions that the retained text needs (\texttt{\textbackslash C}, \texttt{\textbackslash E}, \texttt{\textbackslash R}), with extra wrapper packages (bm, bbm, dsfont, xspace, cancel, mathtools). The delivered numbering is the unit's own. Assets: no retained span contains a figure environment, a graphics-inclusion command, a PDF-page inclusion, a table or a TikZ picture; no image file is copied into this package, the package declares no asset dependency and no image file is redistributed with it. Licence: version arXiv:2606.01832v1 of this article is distributed under the Creative Commons Attribution 4.0 International licence, as recorded in the arXiv licence metadata for this exact version and shown on the abstract page https://arxiv.org/abs/2606.01832v1. A full rights notice is delivered at licenses/arxiv-2606.01832-CC-BY-4.0.txt. Characters: every retained excerpt is pure ASCII, so the delivered engine, XeLaTeX with its default Latin Modern text font, reports no missing character.
\begin{Lemma}\label{asymptotic_powerseries}
    Let $(p_n)$ be a sequence of real numbers. Let $z_0>0$. Assume that
\begin{equation}
\sum_{n=0}^\infty p_n z^n = \frac{f(z)}{z_0-z},
\end{equation}
for complex $|z| < z_0$. Further assume that $f$ is analytic on $|z|<z_1$ with $z_0<z_1$ and that $f(z_0)\neq 0$. Then
$$
p_n \sim f(z_0) \cdot z_0^{-(n+1)},\qquad \text{as $n\to\infty$.}
$$
\end{Lemma}

\begin{align}\label{root_blue}
    E\left(\theta,\frac{-z}{a+1}\right)=0.
\end{align}

\begin{Lemma}\label{asymptotic_blau_erste}
In the blue region, as $n\to\infty$, it holds for $y\in[-a,1]$ that
    \begin{align*}
       p_n^{a,y}(\theta)\sim \frac{a+1}{E\left(\theta,\frac{-\theta }{\lambda(a+1)}\right)}\, \lambda^{n+1}E\left(\theta,\frac{-\theta y}{\lambda(a+1)}\right),
    \end{align*}
    where $\lambda=\lambda(\theta,a)$ is given as the smallest real positive root equation of \eqref{root_blue}.
\end{Lemma}

\begin{Lemma}\label{generatingfunction_gelbe_region}
    In the yellow region, on the common domain of convergence, we have
    \begin{align*}
        \hat{P}_{a,y,\theta}(z)
        &=\begin{cases}
           \frac{E\left(\theta,\frac{-\theta zy}{a+1}\right)}{E\left(\theta,\frac{za}{\theta(a+1)}\right)-z\cdot\frac{ \theta +a}{\theta(a+1)}E\left(\theta,\frac{za}{a+1}\right)},~~&~-a\leq y \leq \frac{-a}{\theta},
            \\\frac{E\left(\theta,\frac{za}{a+1}\right)}{E\left(\theta,\frac{za}{\theta(a+1)}\right)-z\cdot\frac{\theta +a}{\theta(a+1)}E\left(\theta,\frac{za}{a+1}\right)},~~&~\frac{-a}{\theta}<y\leq1.
        \end{cases}    
    \end{align*}
    % with radius of convergence $|z|<\frac{1}{\lambda}$ with $\lambda$ defined in \eqref{root_yellow_2}.
\end{Lemma}

\begin{equation}\label{root_yellow_2}
\frac{E\left(\theta,\frac{za}{\theta(a+1)}\right)}{E\left(\theta,\frac{za}{a+1}\right)} = z \,\frac{\theta+a}{\theta(a+1)},\qquad \text{for $z\in\C$}.
\end{equation}

\begin{Lemma}\label{lemma_root_yellow}
We have $z_0\in\R$, $z_0>0$, for any other solution $z_1\in\C$ of (\ref{root_yellow_2}) we have $|z_1|>z_0$, and $E(\theta,\frac{za}{a+1})\neq 0$ for all $|z|\leq z_0$.
\end{Lemma}

\begin{Lemma}\label{asymptotic_gelb_eigenfunction}
    In the yellow region, it holds with $\lambda=\lambda(\theta,a)$ as given in \eqref{root_yellow_2} that
    \begin{align*}
       p_n^{a,y}(\theta)\sim c(\theta,a)E\left(\theta,\frac{-\theta y}{\lambda(a+1)}\right)\lambda^{n+1}
   \end{align*}
   for $-a\leq y \leq \frac{-a}{\theta}$; and 
   \begin{align*}
       p_n^{a,y}(\theta)\sim c(\theta,a)E\left(\theta,\frac{a}{\lambda(a+1)}\right)\lambda^{n+1}
   \end{align*}
   for $\frac{-a}{\theta}<y\leq1$, with the abbreviation
    \begin{align*}
       c(\theta, a)\coloneqq\frac{1}{\frac{-a}{\theta(a+1)}\left(E\left(\theta,\frac{a}{\lambda(a+1)}\right)\frac{-\theta}{a}- \frac{\theta +a}{\lambda(a+1)}
        E\left(\theta,\frac{\theta a}{\lambda(a+1)}\right)\right)} % \label{def_c1_yellow}.
   \end{align*}
   \end{Lemma}

\begin{proof}
Similarly to the proof of Lemma~\ref{asymptotic_blau_erste}, we construct a function $f(z)$ with
    \begin{align}\label{gelb_f}
        \sum_{n=0}^{\infty}p_n^{a,y}(\theta)z^n=\frac{f(z)}{z_0-z},~~~~\text{for all}~|z|<z_0,
    \end{align}
   which is analytic on $|z|<z_1$ with some $z_0<z_1$ and $f\! \left(z_0\right)\neq 0$. Applying Lemma~\ref{asymptotic_powerseries} will then yield the claim.
\\Recall that the generating function in Lemma~\ref{generatingfunction_gelbe_region} has the structure 
\begin{align*}
    \hat{P}_{a,y,\theta}(z)=\frac{u(z)}{g(z)}
\end{align*}
with $g(z)\coloneqq E\left(\theta,\frac{za}{\theta(a+1)}\right)-z\frac{\theta +a}{\theta(a+1)}E\left(\theta,\frac{za}{a+1}\right)$ and $u(z)$ depending on the choice of $y$: 
\begin{align*}
    u(z)\coloneqq\begin{cases}
         E\left(\theta,\frac{-\theta zy}{a+1}\right),~~~&-a\leq y \leq \frac{-a}{\theta},
         \\E\left(\theta,\frac{za}{a+1}\right),~~&\frac{-a}{\theta}<y\leq1.
    \end{cases}
\end{align*}
Moreover, $z_0>0$ is the root with the smallest absolute value of $g$. Set $v(z)\coloneqq g(z)\frac{1}{z_0-z}$.
We show that the function $f$ defined by $f(z)\coloneqq\frac{u(z)}{v(z)}$ satisfies the assumptions stated above.

By Lemma~\ref{lemma_root_yellow}, $z_0$ is the unique zero of $g$ of smallest modulus with all other zeros lying outside $|z|<z_1$ for some $z_1>z_0$. We will see below that $g'(z_0)<0$, implying that $z_0$ is a simple zero. Therefore, $v$ is analytic and does not vanish on a disk $|z|<z_1$ with $z_1>z_0$. Further, $u$ is entire so that $f$ is analytic on $|z|<z_1$. It remains to be seen that $f(z_0)\neq 0$ in order to apply Lemma~\ref{asymptotic_powerseries} and to compute $f(z_0)$.

For this purpose, note that $v(z_0)=-g'(z_0)$, which can be computed readily:
\begin{align*}
    g'(z_0) & = E\left(\theta,\frac{z_0 a}{a+1}\right) \frac{a}{\theta(a+1)} - \frac{\theta +a}{\theta(a+1)}E\left(\theta,\frac{z_0 a}{a+1}\right) - z_0 \frac{\theta +a}{\theta(a+1)}E\left(\theta,\frac{\theta z_0a}{a+1}\right) \frac{a}{a+1} 
    \\
    & = \frac{-1}{a+1} E\left(\theta,\frac{z_0 a}{a+1}\right)  - z_0 \frac{a(\theta +a)}{\theta(a+1)^2}E\left(\theta,\frac{\theta z_0a}{a+1}\right) < 0,
\end{align*}
because the $E$-terms are positive (due to Lemma~\ref{lemma_root_yellow}) while the prefactors are negative since $\theta\in[-1,0)$ and $a<-\theta$.

It remains to show $u(z_0)\neq 0$ with the help of Lemma~\ref{lemma_root_yellow}. For  $-a\leq y \leq \frac{-a}{\theta}$, we have $u(z)= E\left(\theta,\frac{-\theta zy}{a+1}\right)$ and $|\theta yz_0|\le z_0 a$ and thus,
%\begin{align*}
$u(z_0)= E\left(\theta,\frac{-\theta z_0y}{a+1}\right)\neq 0$. %\end{align*}
  For $\frac{-a}{\theta}<y\leq1$, 
  $u(z)= E\left(\theta,\frac{za}{a+1}\right)$ and
 %  \begin{align*}
      $u(z_0)= E\left(\theta,\frac{z_0a}{a+1}\right)%\overset{}{=}\frac{E\left(\theta,\frac{z_0a}{\theta(a+1)}\right)\theta(a+1)}{z_0(\theta a+a)}
      \neq 0$.
  % \end{align*}
   This gives us $f(z_0)\neq 0$ as well as the value of $f(z_0)=f(\frac{1}{\lambda})$ in both of the cases and hence Lemma~\ref{asymptotic_powerseries} yields the claims. 
% By construction, \eqref{gelb_f} holds.
%     $u$ and $g$ are analytic and $v$ is analytic on $|z|<z_0$, so for $f$ to be analytic on $|z|<z_1$ for some $z_1>z_0$, it suffices to show that $v$ has no zero or singularity in $z_0$. This is done by analyzing the complex Taylor series expansion of $g$ in $z_0$. To get a shorter representation of the Taylor series, we compute the derivatives $g^{(n)}(z_0)$ first. With $g_1(z)\coloneqq E\left(\theta,\frac{za}{\theta(a+1)}\right)$ and  $g_2(z)\coloneqq z\cdot\frac{\theta +a}{\theta(a+1)}E\left(\theta,\frac{za}{a+1}\right)$, we have $g=g_1-g_2$. The differential equation for the deformed exponential function (\ref{eqn:dgl_n_exponential}) gives us 
%     \begin{align*}
%         g_1^{(n)}(z)
%         &=\theta^{n(n-1)/2}\left(\frac{a}{\theta(a+1)}\right)^nE\left(\theta,\frac{\theta^nza}{\theta(a+1)}\right)
%         =\theta^{n(n-1)/2-n}\left(\frac{a}{a+1}\right)^nE\left(\theta,\frac{\theta^{n-1}za}{a+1}\right)
%     \end{align*}
%     as well as
%     \begin{align*}
%         g_2^{(n)}(z)
%         &=\frac{\theta +a}{\theta(a+1)}\sum_{k=0}^n\binom{n}{k}\left(\frac{\partial^k }{\partial z^k } z\right) 
%         \left(\frac{\partial^{n-k} }{\partial z^{n-k} }E\left(\theta,\frac{za}{a+1}\right)\right)
%      %   \\&=\frac{\theta +a}{\theta(a+1)}\sum_{k=0}^n\binom{n}{k}\left(\frac{\partial^k }{\partial z^k } z\right) 
%      %   \theta^{({n-k})((n-k)-1)/2}\left(\frac{a}{a+1}\right)^{n-k}E\left(\theta,\frac{\theta^{n-k}za}{a+1}\right)
%         \\&=\frac{\theta +a}{\theta(a+1)}\bigg(\binom{n}{0}\left(\frac{\partial^0 }{\partial z^0} z\right) 
%         \theta^{{n}(n-1)/2}\left(\frac{a}{a+1}\right)^{n}E\left(\theta,\frac{\theta^{n}za}{a+1}\right)
%         \\&~~~~~~~~~~~~~~~~~~~~~~~+\binom{n}{1}\left(\frac{\partial^1 }{\partial z^1} z\right) 
%         \theta^{({n-1})(n-2)/2}\left(\frac{a}{a+1}\right)^{n-1}E\left(\theta,\frac{\theta^{n-1}za}{a+1}\right)\bigg)
%         \\&=\frac{\theta +a}{\theta(a+1)}\bigg(z 
%         \theta^{n(n-1)/2}\left(\frac{a}{a+1}\right)^{n}E\left(\theta,\frac{\theta^{n}za}{a+1}\right)\\&~~~~~~~~~~~~~~~~~~~~~~~+ n 
%         \theta^{n({n-1})/2}\theta^{-n+1}\left(\frac{a}{a+1}\right)^{n-1}E\left(\theta,\frac{\theta^{n-1}za}{a+1}\right)\bigg).
%     \end{align*}
%     Summing up, this gives 
%     \begin{align*}
%         g^{(n)}(z_0)
%         &=g_1^{(n)}(z_0)-g_2^{(n)}(z_0)
%         %\\&=\theta^{n(n-1)/2-n}\left(\frac{a}{a+1}\right)^nE\left(\theta,\frac{\theta^{n-1}z_0a}{a+1}\right)
%         %\\&~~~~~-\frac{\theta +a}{\theta(a+1)}z_0 \theta^{n(n-1)/2}\left(\frac{a}{a+1}\right)^{n}E\left(\theta,\frac{\theta^{n}z_0a}{a+1}\right) 
%       % \\&~~~~~- \frac{\theta +a}{a+1}n \theta^{n(n-1)/2-n+1}\left(\frac{a}{a+1}\right)^{n-1}E\left(\theta,\frac{\theta^{n-1}z_0a}{a+1}\right)
%         \\&=\theta^{n(n-1)/2-n}\left(\frac{a}{a+1}\right)^nE\left(\theta,\frac{\theta^{n-1}z_0a}{a+1}\right)\left(1-\frac{n(\theta+a)}{a}\right)
%         \\&~~~~~-\frac{\theta +a}{\theta(a+1)}z_0 
%         \theta^{n(n-1)/2}\left(\frac{a}{a+1}\right)^{n}E\left(\theta,\frac{\theta^{n}z_0a}{a+1}\right) .
%         %\\&=\theta^{n(n-1)/2-n}\left(\frac{a}{a+1}\right)^nE\left(\theta,\frac{\theta^{n-1}z_0a}{a+1}\right)\left(1-\frac{n(\theta+a)}{a}\right)
%         %\\&~~~~~-\frac{E\left(\theta,\frac{z_0a}{\theta(a+1)}\right)}{E\left(\theta,\frac{z_0a}{a+1}\right)}
%         %\theta^{n(n-1)/2}\left(\frac{a}{a+1}\right)^{n}E\left(\theta,\frac{\theta^{n}z_0a}{a+1}\right),
%        % \\&=\frac{\theta^{(n-1)(n-2)/2}}{\theta}\left(\frac{a}{a+1}\right)^nE\left(\theta,\frac{\theta^{n-1}z_0a}{a+1}\right)\left(1-n(\theta+1)\right)
%         %\\&~~~~~-\frac{\theta^{n(n-1)/2}}{\theta}z_0 (\theta +1)
%         %\left(\frac{a}{a+1}\right)^{n+1}E\left(\theta,\frac{\theta^{n}z_0a}{a+1}\right)  
%     \end{align*}
%     By definition of $z_0$, $g(z_0)=0$.
%     Thus, for the Taylor series expansion, it holds
%     \begin{align*}
%        Tg(z;z_0)&=\sum_{n=0}^\infty \frac{g^{(n)}(z_0)}{n!}(z-z_0)^n=\sum_{n=1}^\infty \frac{g^{(n)}(z_0)}{n!}(z-z_0)^n.
%        %\\&=\sum_{n=0}^\infty \frac{\frac{\theta^{(n-1)(n-2)/2}}{\theta}\left(\frac{a}{a+1}\right)^nE\left(\theta,%\frac{\theta^{n-1}z_0a}{a+1}\right)\left(1-n(\theta+1)\right)}{n!}(z-z_0)^n
%        %\\&~~~~~-\sum_{n=0}^\infty \frac{\frac{E\left(\theta,\frac{z_0a}{\theta(a+1)}\right)}{E\left(\theta,\frac{z_0a}{a+1}\right)}
%       %  \theta^{n(n-1)/2}\left(\frac{a}{a+1}\right)^{n}E\left(\theta,\frac{\theta^{n}z_0a}{a+1}\right)}{n!}(z-z_0)^n.
%    \end{align*}
%    By definition of $v$, this gives
%    \begin{align*}
%        v(z)&=\frac{1}{z_0-z}\sum_{n=1}^\infty \frac{g^{(n)}(z_0)}{n!}(z-z_0)^n.
%        \\&=-\sum_{n=1}^\infty \bigg({\theta^{-n}E\left(\theta,\frac{\theta^{n-1}z_0a}{a+1}\right)\left(1-\frac{n(\theta+a)}{a}\right)}
%        - \frac{\theta +a}{\theta(a+1)}z_0
%         E\left(\theta,\frac{\theta^{n}z_0a}{a+1}\right)\bigg)\\&~~~~~~~~~~~~~~~~~~~~~\cdot\frac{\theta^{n(n-1)/2}}{n!}\left(\frac{a}{a+1}\right)^n(z-z_0)^{n-1},
%    \end{align*}
%   which has no singularity in $z=z_0$. Furthermore,  for $z=z_0$, only the first term of the sum is not equal to zero, so
%   \begin{align*}
%       v(z_0)
%      % &=\frac{-1}{\theta}\frac{a}{a+1}E\left(\theta,\frac{z_0a}{a+1}\right)\frac{-\theta}{a}+\frac{E\left(\theta,\frac{z_0a}{\theta(a+1)}\right)}{E\left(\theta,\frac{z_0a}{a+1}\right)}
%       %  \frac{a}{a+1}E\left(\theta,\frac{\theta z_0a}{a+1}\right)
%     &=\frac{-a}{\theta(a+1)}\left(E\left(\theta,\frac{z_0a}{a+1}\right)\frac{-\theta}{a}- z_0\frac{\theta +a}{a+1}
%         E\left(\theta,\frac{\theta z_0a}{a+1}\right)\right),
%   \end{align*}
%   which is non-zero: The first factor is positive since we excluded $a=0$ and the second term is positive because $a<-\theta$ and $z_0>0$, so it is the difference of two terms with different signs, which we already justified in Lemma~\ref{lemma_root_yellow}. So we have shown that $f$ is analytic on $|z|<z_1$ for some $z_0<z_1$. The last step is to show that $f(z_0)\neq 0$. We already computed $v(z_0)$, so it is left to show $u(z_0)\neq 0$ with the help of Lemma~\ref{lemma_root_yellow}. For  $-a\leq y \leq \frac{-a}{\theta}$, we have $u(z)= E\left(\theta,\frac{-\theta zy}{a+1}\right)$ and $|\theta yz_0|\le z_0 a$ and thus,
% \begin{align*}
%     u(z_0)= E\left(\theta,\frac{-\theta z_0y}{a+1}\right)\neq 0.
% \end{align*}
%   For $\frac{-a}{\theta}<y\leq1$, it is
%   $u(z)= E\left(\theta,\frac{za}{a+1}\right)$, so 
%   \begin{align*}
%       u(z_0)= E\left(\theta,\frac{z_0a}{a+1}\right)%\overset{}{=}\frac{E\left(\theta,\frac{z_0a}{\theta(a+1)}\right)\theta(a+1)}{z_0(\theta a+a)}
%       \neq 0.
%   \end{align*}
%    This gives us $f(z_0)\neq 0$ in both of the cases and thus Lemma~\ref{asymptotic_powerseries} yields the claims. 
%   %  Therefore, we have for $-a\leq y \leq \frac{-a}{\theta}:$
%   % \begin{align*}
%   %      f(z_0)&
%   %      =\frac{E\left(\theta,\frac{-\theta z_0y}{a+1}\right)}{\frac{-a}{\theta(a+1)}\left(E\left(\theta,\frac{z_0a}{a+1}\right)\frac{-\theta}{a}- z_0\frac{\theta +a}{a+1}
%   %       E\left(\theta,\frac{\theta z_0a}{a+1}\right)\right)}\neq 0
%   %   \end{align*}
%   %   and for $\frac{-a}{\theta}<y\leq1$,
%   %   \begin{align*}
%   %       f(z_0)&
%   %       =\frac{E\left(\theta,\frac{z_0a}{a+1}\right)}{\frac{-a}{\theta(a+1)}\left(E\left(\theta,\frac{z_0a}{a+1}\right)\frac{-\theta}{a}- z_0\frac{\theta +a}{a+1}
%   %       E\left(\theta,\frac{\theta z_0a}{a+1}\right)\right)}\neq 0.
%   %  \end{align*}
%   %  To abbreviate the notation, we use $c(\theta,a)$ defined in \eqref{def_c1_yellow}. Then Lemma~\ref{asymptotic_powerseries} yields the claims. 
%    % with $z_0=\frac{1}{\lambda}:$
%    % \begin{align*}
%    %     p_n^{a,y}(\theta)\sim f\left(\frac{1}{\lambda}\right)\lambda^{n+1}=c(\theta,a)E\left(\theta,\frac{-\theta y}{\lambda(a+1)}\right)\lambda^{n+1}
%    % \end{align*}
%    % for $-a\leq y \leq \frac{-a}{\theta}$ and 
%    % \begin{align*}
%    %     p_n^{a,y}(\theta)\sim f\left(\frac{1}{\lambda}\right)\lambda^{n+1}=c(\theta,a)E\left(\theta,\frac{a}{\lambda(a+1)}\right)\lambda^{n+1}
%    % \end{align*}
%    % for $\frac{-a}{\theta}<y\leq1$.
   \end{proof}

\end{document}
